\label{sekcja_piaty_postulat}

Tak jak pisaliśmy na stronie \pageref{pageref:fifth_postulate}, piąty postulat Euklidesa będzie swoją długością budzić podejrzenia; wiele osób spróbuje usunąć go z listy przez wyprowadzenie z pozostałych albo dodanie innego, prostszego aksjomatu.
Za każdym razem okaże się, że w takim uzasadnieniu przemycone będzie zdanie równoważne piątemu postulatowi.
Historię przedstawią Hartshorne \cite[s. 296-304]{hartshorne_2000} i Eves \cite[s. 283-288]{eves1_1972}; my uzupełniamy ją o drobne, brakujące szczegóły.
Opowiemy o Proklosie, Claviusie, Wallisie, Clairaulcie, Simsonie, Playfairze i Bolyai.
% Między innymi Ptolemeusz, Proklos, Ibn al-Hajsam, Omar Chajjam, Nasir ad-Din at-Tusi oraz jego syn Sadr ad-Din, Christopher Clavius: jezuita z Collegio Romano, Giordano Vitale, John Wallis w wykładzie z 1663 roku do którego zainspiruje go traktat Sadr ad-Dina.
% Każdy z nich zostawi po sobie zdanie, które zastępuje piąty postulat -- o nich opowiemy w następnej sekcji. % lang-pl
Innych (nieopisanych tutaj) prób podejmą się B.~F. Thibaut (1809), J.~D. Gergonne (1812) i J.~K.~F. Hauff (1819); opisy tych prób znajdziemy u Evesa \cite[s. 289]{eves1_1972} oraz Hartshorne'a \cite[s. 304]{hartshorne_2000}. % lang-pl

Potrzebujemy wtrącić treść aksjomatu Archimedesa, wbrew nazwie nadanej przez Ottona Stolza w~1883 roku sformułowanego po raz pierwszy przez Eudoksosa z Knidos, a nie Archimedesa.
Głosi, że każdy odcinek jest krótszy od pewnej wielokrotności dowolnego innego odcinka.
Wynika stąd:

\begin{proposition}
[aksjomat Arystotelesa] % lang-pl
	Dla każdego kąta ostrego $\angle XOY$ i odcinka $AB$ istnieje punkt $Q$ na półprostej $OY$ taki, że odcinek $PQ$ (gdzie $P$ jest rzutem $Q$ na półprostą $OX$) jest dłuższy od odcinka $AB$.
\index{aksjomat!Arystotelesa}%
\end{proposition}

Ten aksjomat znajdzie się w dziele Arystotelesa \emph{,,Περὶ οὐρανοῦ''} (Peri urano, dosłownie: o niebie) z 350 roku przed Chrystusem.
% On the Heavens (Greek: ; Latin: De Caelo or De Caelo et Mundo) is Aristotle's chief cosmological treatise: written in 350 BCE,[1] it contains his astronomical theory and his ideas on the concrete workings of the terrestrial world. It should not be confused with the spurious work On the Universe (De Mundo, also known as On the Cosmos).
Jego implikacje opisze Hartshorne \cite[s. 319-326]{hartshorne_2000}.

Hartshorne \cite[s. 297]{hartshorne_2000} pisze, że to wystarczyło do wyprowadzenia pewnej formy aksjomatu Playfaira, o którym później:

\begin{proposition}
[lemat Proklosa] % lang-pl
	Prosta, która przecina jedną z dwóch równoległych prostych, przecina także tę drugą.
\index{lemat!Proklusa}%
\end{proposition}

Zawiera błąd: Proklos mówi o~odległości między dwiema prostymi tak, jakby każdy punkt jednej leżał w takiej samej odległości od drugiej.
Ale definicja równoległych mówi jedynie, że te się nie przecinają; to, co dołożył od siebie Proklos, w połączeniu z aksjomatem Arystotelesa jest równoważne piątemu postulatowi.

Rozumowanie Proklosa musiało być jednak akceptowane przez ówczesnych, bo przytoczy je bez krytycznych uwag F.~Commandino w swoim wydaniu Euklidesa z 1575 roku.
Ten sam błąd popełnią później Jacques Peletier (w 1557 roku), Andrea Tacquet (w 1654 roku) i wiele innych, niewymienionych tu osób.
Tacquet korzysta z:

\begin{proposition}
[aksjomat Claviusa] % lang-pl
	Zbiór punktów leżących po jednej stronie danej prostej i jednakowo od niej odległych jest prostą. % lang-pl
\index{aksjomat!Claviusa}%
\end{proposition}

Ale aksjomat Claviusa razem z aksjomatem Archimedesa implikują piąty postulat; istnieje też niearchimedesowska płaszczyzna, na której aksjomat Claviusa jest spełniony, ale piąty postulat nie (ćwiczenie 33.7 u Hartshorne'a \cite[s. 302]{hartshorne_2000}).

Kiedyś między 1616 oraz 1703 rokiem, John Wallis poda

\begin{proposition}
[aksjomat Wallisa] % lang-pl
	Dla każdego odcinka $DE$ i trójkąta $\triangle ABC$ istnieje trójkąt do niego podobny, $\triangle A'B'C'$, którego bok $A'B'$ jest dłuższy od $DE$. % lang-pl
\index{aksjomat!Wallisa}%
\end{proposition}

Aksjomat Wallisa jest nieprawdziwy w geometrii niearchimedesowej, gdzie istnieją trójkąty podobne o różnych rozmiarach.
Wynika z niego aksjomat Playfaira.

Alexis Claude de Clairault, matematyk\footnote{Znany również z różniczkowego równania $y = x y' + f(y')$.} francuski, zechce uczynić geometrię bardziej przystępną dla uczniów.
Będzie krytykować nauczanie oparte na definicjach i skomplikowanych dowodach, zamiast tego stawiając na przykłady praktyczne, takie jak pomiary terenu.
W \emph{,,Elémens de géométrie''} z 1741 roku rozwinie teorię prostych równoległych z konstrukcji prostokąta o zadanej podstawie.
Dlatego jest sens mówić o:

\begin{proposition}
[aksjomat Clairaulta] % lang-pl
	Niech $AC$ i $BD$ będą równymi odcinkami prostopadłymi do odcinka $AB$. % lang-pl
	Wtedy kąty przy $C$ i $D$ są proste, czyli $ABCD$ jest prostokątem. % lang-pl
\index{aksjomat!Clairauta}%
\end{proposition}

Trochę później Robert Simson wyda jedną z ważniejszych edycji ,,Elementów'' Euklidesa. % 1756
\index[persons]{Simson, Robert}%
Postanowi usunąć zniekształcenia i~błędy wprowadzone przez wcześniejszych redaktorów, celem przywrócenia dziełu ścisłej formy.
Uwierzy, że wszystko, co matematycznie poprawne i rygorystyczne, musiało pochodzić od samego Euklidesa.
Na potrzeby dowodu piątego postulatu wprowadzi

\begin{proposition}
[aksjomat Simsona] % lang-pl
	Prosta nie może najpierw zbliżyć się do innej prostej, a potem oddalić od niej, zanim ją przetnie; podobnie nie może najpierw oddalić się od innej prostej, a potem do niej zbliżyć; nie może też przez jakiś czas zachowywać stałą odległość od innej prostej, a potem się do niej zbliżyć lub od niej oddalić. % lang-pl
\index{aksjomat!Simsona}%
\end{proposition}

Korzystając z tego aksjomatu oraz (niejawnie) z aksjomatu Archimedesa, poprawnie wyprowadzi pięć twierdzeń, z których ostatnie będzie piątym postulatem.
(Pierwsze wydanie Simsona pojawi się w 1756 roku).
Aksjomaty Claviusa, Clairaulta, Simsona są sobie równoważne (patrz Hartshorne \cite[s. 302]{hartshorne_2000}).

\begin{proposition}
[aksjomat Playfaira] % lang-pl
	Dwie przecinające się proste nie mogą być równoległe do tej samej prostej
\index{aksjomat!Playfaira}%
\end{proposition}

Równoważnie, przez punkt poza prostą można przeprowadzić co najwyżej jedną prostą równoległą do wyjściowej.
John Playfair, autor powyższego zdania, umieści je w swoim wydaniu Elementów z~1795 roku z komentarzem, że znał wydanie Simsona.

Wreszcie Farkas Bolyai, ojciec Jánosa, zaproponuje:

\begin{proposition}
[aksjomat Bolyaia]
	Każde trzy niewspółliniowe punkty leżą na jednym okręgu.
\end{proposition}

Implikuje to piąty postulat, jak każe wykazać Hartshorne w ćwiczeniu.


Bez aksjomatu Archimedesa, piąty postulat jest mocniejszym stwierdzeniem od tego, że istnieje prostokąt (Dehn, 1900).
% TODO: Dehn, Max (1900), "Die Legendre'schen Sätze über die Winkelsumme im Dreieck", Mathematische Annalen, 53 (3): 404–439, doi:10.1007/BF01448980, S2CID 122651688
To zaś jest mocniejsze niż aksjomat wprowadzony i zbadany przez Friedrich Bachmann w 1964 roku, w artykule \emph{Zur Parallelenfrage}: % lang-pl
\index[persons]{Bachmann, Friedrich} % lang-pl

\begin{axiom} % lang-pl
[Lotschnittaxiom] % lang-pl
\label{lotschnittaxiom}%
\index{aksjomat!przecinających się prostopadłych (Lotschnittaxiom)} % lang-pl
	Prostopadłe poprowadzone z obu ramion kąta prostego przecinają się. % lang-pl
\end{axiom} % lang-pl

Lotschnittaxiom jest tak samo mocny jak aksjomat Legendre dla kąta prostego albo fakt, że każde trzy równoległe proste można przeciąć czwartą.
Pambuccian pokaże w 1994 roku, że w obecności aksjomatu Archimedesa, aksjomat \ref{lotschnittaxiom} jest równoważny piątemu postulatowi.

% Źródło: https://en.wikipedia.org/wiki/Lotschnittaxiom
% Bachmann, Friedrich (1964), "Zur Parallelenfrage", Abhandlungen aus dem Mathematischen Seminar der Universität Hamburg, 27 (3–4): 173–192, doi:10.1007/BF02993215, S2CID 186240918.
% Pambuccian, Victor (1994), "Zum Stufenaufbau des Parallelenaxioms", Journal of Geometry, 51 (1–2): 79–88, doi:10.1007/BF01226859, hdl:2027.42/43033, S2CID 28056805